\begin{afterword}
\summaryhead

The post starts from a reader's definition of art, ``that which is designed for the purpose of creating a reaction in an audience.'' It replies that a word has no meaning floating in the void, waiting to be discovered by finding the right definition. The real task is to find which things cluster together and which share a common cause: to carve reality at its joints. A word ``eluctromugnetism'' that took in lightning, compasses and animal magnetism and left out light would be a wrong word.

Given a list of things called art and not art, the author proposes a rule: deliberate effort meant to produce aesthetic emotion. Someone could dispute the rule, or dispute the list itself, on the presumption that some underlying pattern generates it. Animal magnetism could be doubted before anyone knew what electricity and magnetism share, and a list of fish that includes dolphins is wrong. So both kinds of definition, rules and lists, can be wrong, and one should admit that one's definitions are guesses that can be mistaken.

\responsehead

The post's main claim is stated plainly and is sound. A category is good or bad according to whether it groups things that are really alike or share a cause, and a definition, whether a rule or a list, can fail that test. The two scientific cases support it. Animal magnetism was set aside by a royal commission in 1784, decades before anyone connected electricity and magnetism, and Linnaeus moved whales out of the fishes in 1758.

The opening example does not show what the post uses it to show. The post presents the commenter as someone who has ``pondered the meaning of the word'' and so is looking at the problem the wrong way. Those words are the post's. The comment it links to offered a definition and said nothing about the meaning of the word. The definition offered is the same kind of thing as the author's own a few paragraphs later: a statement of what the things called art have in common, in terms of intent and the audience's response. The post treats its own guess as a legitimate attempt to describe a cluster. Apart from the words it added, it gives no reason to treat the commenter's guess as a search for an essence.

The art example is also where the main claim has the least support. The post holds that a list can be disputed on ``the presumption'' that some pattern generates it. For the invented word ``eluctromugnetism'' and for fish it names what settles the dispute: a later finding about what the members share. For art it names no finding that would decide whether Modern Art belongs, and it closes the question with a joke about philistines. The art case shows that people can dispute where a boundary goes. Only the scientific cases support the claim that a boundary can be wrong.

In short: a sound claim that definitions can fail to ``carve reality at its joints,'' in the post's phrase, by grouping things that are not really alike, supported by its scientific cases, opened with a commenter to whom the post attributes, in words of its own, the error it criticizes, and illustrated with an art example that shows a dispute without showing how it could be settled.
\end{afterword}
